\section{Тестирование и бенчмарк}

Корректность программы проверялась на примере из условия, на граничных случаях и на случайных малых
тестах. Для случайных тестов результат сравнивался с прямым перебором всех позиций текста и всех
образцов. Отдельно проверялись следующие ситуации:
\begin{itemize}
\item образец длины \(1\);
\item несколько образцов, заканчивающихся в одной позиции текста;
\item образцы, являющиеся суффиксами других образцов;
\item числа с ведущими нулями во входе;
\item значение \(2^{32} - 1\);
\item пустые строки в текстовой части входа.
\end{itemize}

Для оценки производительности был подготовлен отдельный бенчмарк. Данные генерируются скриптом
\texttt{gen\_benchmark\_data.py}, а измерения выполняются программой \texttt{benchmark.cpp}. Генерация
входа и разбор данных выполняются вне измеряемого участка. Время замеряется отдельно для построения
автомата и для прохода по уже сохранённому в памяти тексту. Для каждого случая выполнено 5 запусков,
в таблицах приведены медианные значения. Полные результаты сохранены в файле
\texttt{results/benchmark.csv}.

Команды запуска:
\begin{alltt}
$ g++ -std=c++14 -O2 -Wall -Wextra -pedantic main.cpp -o lab4.out
$ bash run_benchmarks.sh
\end{alltt}

\subsection{Сравнение с наивным поиском}

На малых и средних тестах алгоритм Ахо--Корасик сравнивался с прямым перебором всех позиций текста и
всех образцов. В последнем столбце указано отношение \texttt{Naive / AC}; значение больше 1 означает,
что автомат работает быстрее. Для всех строк число найденных вхождений у двух алгоритмов совпало.

\begin{center}
\begin{tabular}{|r|r|r|r|r|r|}
\hline
\rowcolor{lightgray}
Образцов & Длина текста & AC build, \(\mu s\) & AC search, \(\mu s\) & Naive, \(\mu s\) & Naive / AC \\
\hline
10 & 20\,000 & 7 & 189 & 238 & \(1.21\times\) \\
\hline
50 & 20\,000 & 11 & 110 & 543 & \(4.49\times\) \\
\hline
100 & 20\,000 & 19 & 93 & 1128 & \(10.16\times\) \\
\hline
\end{tabular}
\end{center}

Даже на небольшом тексте преимущество автомата быстро растёт с увеличением числа образцов. Это
ожидаемо: наивный алгоритм заново сравнивает каждый образец с каждой позицией текста, а
Ахо--Корасик после построения автомата выполняет один общий проход.

\subsection{Масштабирование по длине текста}

В следующей серии число образцов фиксировано и равно \(500\), длина каждого образца равна \(5\).
Менялась только длина текста. Наивный алгоритм здесь не запускался, так как его стоимость уже
существенно выше и не нужна для проверки масштабирования автомата.

\begin{center}
\begin{tabular}{|r|r|r|r|r|}
\hline
\rowcolor{lightgray}
Чисел в тексте & Вхождений & Build, \(\mu s\) & Search, \(\mu s\) & Total, \(\mu s\) \\
\hline
50\,000 & 50 & 99 & 266 & 363 \\
\hline
100\,000 & 100 & 100 & 534 & 633 \\
\hline
500\,000 & 500 & 103 & 2686 & 2794 \\
\hline
1\,000\,000 & 1000 & 102 & 5455 & 5575 \\
\hline
\end{tabular}
\end{center}

Время построения практически не меняется, потому что набор образцов один и тот же. Время поиска растёт
почти линейно относительно числа чисел в тексте: при увеличении текста с \(100000\) до \(1000000\)
время прохода выросло примерно в \(10.2\) раза.

\subsection{Масштабирование по числу образцов}

В этой серии длина текста фиксирована и равна \(200000\) чисел, длина каждого образца равна \(5\).
Менялось число образцов, то есть суммарная длина trie.

\begin{center}
\begin{tabular}{|r|r|r|r|r|}
\hline
\rowcolor{lightgray}
Образцов & Суммарная длина & Build, \(\mu s\) & Search, \(\mu s\) & Total, \(\mu s\) \\
\hline
100 & 500 & 19 & 1292 & 1317 \\
\hline
500 & 2500 & 105 & 1102 & 1202 \\
\hline
1000 & 5000 & 197 & 1102 & 1301 \\
\hline
5000 & 25000 & 984 & 1066 & 2050 \\
\hline
\end{tabular}
\end{center}

Время построения растёт вместе с суммарной длиной образцов, что соответствует оценке \(O(P)\).
Время поиска при этом остаётся одного порядка: текст имеет ту же длину, а сгенерированные переходы
хранятся разреженно в \texttt{unordered\_map}. Различия между строками связаны с константами хеш-таблиц
и конкретным набором переходов.

\subsection{Нагрузка на выходные ссылки}

Отдельно проверялся случай с большим числом найденных ответов. Для этого использовались образцы
\(x\), \(x\,x\), \(\ldots\), \(x^k\), а текст состоял из \(100000\) одинаковых чисел \(x\). В таком
тесте почти в каждой позиции заканчивается сразу несколько образцов.

\begin{center}
\begin{tabular}{|r|r|r|r|}
\hline
\rowcolor{lightgray}
Макс. длина образца & Найдено вхождений & Search, \(\mu s\) & Total, \(\mu s\) \\
\hline
5 & 499\,990 & 470 & 471 \\
\hline
10 & 999\,955 & 821 & 823 \\
\hline
20 & 1\,999\,810 & 2210 & 2216 \\
\hline
\end{tabular}
\end{center}

Этот тест показывает вклад выходных ссылок и количества ответов. При одинаковой длине текста время
увеличивается вместе с числом найденных вхождений, что соответствует оценке поиска \(O(T + Ans)\).
\pagebreak
